\documentclass[11pt]{article}
\usepackage[margin=1in]{geometry}
\usepackage[T1]{fontenc}
\usepackage{lmodern}
\usepackage{booktabs}
\usepackage{tabularx}
\usepackage{longtable}
\usepackage{array}
\usepackage{xcolor}
\usepackage{hyperref}
\usepackage{fvextra}
\setlength{\parindent}{0pt}
\setlength{\parskip}{5pt}
\title{Online Appendix\\Operational AI Disclosure in SEC Filings}
\author{Anonymous manuscript}
\date{}

\begin{document}
\maketitle
\tableofcontents

\section{Stable-lexicon and firm-level checks}

\subsection{Stable-lexicon check}

The benchmark candidate frame used one lexicon frozen before classification
and applied without period-specific changes. As an additional robustness
check, the stable-lexicon specification removes terms specific to the
post-ChatGPT generative-AI vocabulary: \textit{generative AI},
\textit{generative artificial intelligence}, \textit{gen AI}, \textit{large
language model(s)}, \textit{LLM(s)}, \textit{foundation model(s)}, and
\textit{transformer model(s)}. A passage is retained if it contains at least
one other frozen term. Human correction is re-estimated for this restricted
domain with the original probability-audit weights; it is not a model-only
subsample calculation.

\begin{table}[htbp]
\centering
\caption{Stable-lexicon robustness by period}
\small
\begin{tabularx}{\textwidth}{@{}Xrrr@{}}
\toprule
Quantity & 2017--2022 & 2023--2025 & Post minus pre \\
\midrule
Stable-lexicon candidate passages & 6,568 & 18,703 & --- \\
Attention per 100 eligible filings & 29.27 & 175.98 & --- \\
Corrected operational passages per 100 filings
  & 7.55 & 23.53 & --- \\
Corrected operational share
  & 25.79\% & 13.37\% & $-12.42$ pp \\
Approximate 95\% interval for share & [20.14, 31.45] &
  [10.55, 16.19] & [$-18.75$, $-6.10$] pp \\
\bottomrule
\end{tabularx}
\end{table}

The substantive result is unchanged. The benchmark decline is $-12.31$
percentage points; under the stable lexicon it is $-12.42$ points. Removing
the emerging terms reduces post-period attention and corrected operational
intensity, but does not explain the decline in the operational share.

\subsection{Exploratory firm-year proxy}

The firm-year check aggregates frozen passage predictions, not human-corrected
firm-year adoption. The denominator is candidate-observed firm-years: firm-years
with at least one retained passage. The one-evidence proxy requires at least
one model-positive passage. The stricter proxy requires model-positive
passages in at least two distinct SEC accessions in the same firm-year.

\begin{table}[htbp]
\centering
\caption{Model-label firm-year disclosure proxies}
\small
\begin{tabularx}{\textwidth}{@{}Xrrr@{}}
\toprule
Quantity & 2017--2022 & 2023--2025 & Change \\
\midrule
Candidate-observed firm-years & 788 & 983 & --- \\
At least one positive passage & 50.51\% (398) & 37.95\% (373)
  & $-12.56$ pp \\
Positive passages in at least two accessions & 14.47\% (114) &
  11.70\% (115) & $-2.77$ pp \\
\bottomrule
\end{tabularx}
\end{table}

These proxies describe disclosure persistence within the candidate frame.
They are not estimates of the percentage of firms adopting AI: candidate-free
firm-years are outside the denominator, repeated text can represent one
deployment, and passage-level error correction does not transfer mechanically
to a nonlinear at-least-one firm-year indicator.

\section{Frozen retrieval lexicon}

Candidate retrieval used normalized, case-insensitive whole-term matching.
Hyphens, Unicode dash variants, and spaces were normalized for multiword
terms. Any match entered the candidate frame; semantic relevance was decided
downstream. The exact frozen lexicon was:

\VerbatimInput[fontsize=\small]{online_appendix_files/ai_terms_broad.txt}

\section{Frozen LLM prompt}

The classification call supplied the system instruction, allowed labels,
decision rules, output schema, synthetic examples, and a user-message
template shown verbatim below. Three deterministic-seed calls were made per
passage and combined under the frozen aggregation rule.

\VerbatimInput[fontsize=\scriptsize,breaklines=true,
  breakanywhere=true]{online_appendix_files/v4_llm_prompt.json}

\section{Human codebook}

\subsection{Unit and binary target}

Coders labeled the displayed target passage independently using only the
displayed context. The binary reference outcome equals one only when the
status is \texttt{operational}. Operational status requires all four
conditions: (1) the filing firm is the actor; (2) an AI technology is
identified; (3) an action and object are stated; and (4) current or completed
realization is established. Current AI operating inside the firm's product or
service counts, including AI work performed by employees to deliver a service.
A product name, launch, availability, sales statement, demand statement, or
development plan alone does not establish operation.

\subsection{Status codebook}

\begin{longtable}{@{}p{0.20\textwidth}p{0.74\textwidth}@{}}
\toprule
Label & Coding rule \\
\midrule
\endhead
\texttt{operational} & Current or completed AI operation in the filing
firm's own product, service, workflow, or decision. Present-tense operating
capabilities qualify when tied to the firm's current offering. \\
\texttt{pilot\_trial} & Proof of concept, limited test, pilot, beta, or
shadow deployment without evidence of broader current operation. \\
\texttt{planned} & Intention, evaluation, development plan, or future
commitment without established current use. \\
\texttt{irrelevant\_general} & Rhetoric, risk boilerplate, market discussion,
non-AI meanings, product demand without operating capability, or use by a
customer, supplier, competitor, or industry actor outside the filing firm's
offering. \\
\texttt{uncertain} & The displayed evidence does not permit a defensible
choice among the preceding categories. \\
\bottomrule
\end{longtable}

\subsection{Function, outcome, and provider fields}

Functions were multi-label only when the passage explicitly supported more
than one function. The six classes were: core operations and supply chain;
risk, fraud, compliance, and cybersecurity; customer service, sales, and
marketing; product/service development and R\&D; corporate support; and
enterprise-wide/unspecified. AI inside the product sold by the filing firm was
coded as product/service development and R\&D unless another operating
function was explicitly established.

Outcomes were multi-label only when more than one outcome was explicit. The
six classes were: cost/efficiency/speed; revenue/profit margin; quality,
accuracy, or customer experience; risk/loss reduction; innovation, products,
or workforce; and no explicit outcome. Profitability language was assigned
to revenue/profit margin only when the passage explicitly connected AI to
revenue, profit, margin, pricing, monetization, or sales. Customer experience
was coded with quality/accuracy/customer experience, not automatically with
revenue.

The provider-role field used \texttt{yes}, \texttt{no}, or
\texttt{uncertain}. A provider-positive passage could still describe AI
operating inside the filing firm's own product or service; it was not treated
as an external-provider versus internal-user dichotomy. Coders copied the
shortest exact supporting substring as the evidence span and recorded
confidence as high, medium, or low.

\section{Probability-audit sampling groups}

The original 200-item gold audit was drawn after the 2,000-item model pilot.
Its conditional allocation was enriched for model positives and likely
false-negative strata. The final augmentation excluded those 200 exact audited
units, then sampled 800 of the remaining 26,079 candidate passages by simple
random sampling without replacement within prediction-status $\times$ period
$\times$ filing-form $\times$ sector cells. There were 150 realized
augmentation cells. Inclusion probabilities and inverse weights were retained
at the actual-cell level.

\begin{table}[htbp]
\centering
\caption{Audit-stage group sizes}
\small
\begin{tabularx}{\textwidth}{@{}Xrrl@{}}
\toprule
Stage and group & Frame $N$ & Audit $n$ & Design note \\
\midrule
Original pilot: model positive & 223 & 120 &
Conditional draw within the 2,000-item pilot \\
Original pilot: negative boundary & 37 & 21 &
At least one positive model vote \\
Original pilot: negative zero-vote & 1,740 & 59 &
No positive model vote \\
\addlinespace
Augmentation: model positive & 2,929 & 250 &
Allocated across period/form/sector cells \\
Augmentation: negative boundary & 447 & 200 &
Allocated across period/form/sector cells \\
Augmentation: negative zero-vote & 22,703 & 350 &
Allocated across period/form/sector cells \\
\midrule
Final audited total & 26,279 & 1,000 &
200 exact audited units plus 800 probability-audited remainder units \\
\bottomrule
\end{tabularx}
\end{table}

The table reports marginal group sizes for transparency. Analysis used
$\pi_i=n_h/N_h$ and $w_i=1/\pi_i$ from each actual augmentation cell $h$, not
the marginal ratios in the table. The 200 original audited units had analysis
weight one after exclusion from the augmentation frame; their initial
two-stage selection probabilities describe how that development/gold audit
was formed, not their final remainder-cell weight.

\section{Sampling-weighted confusion matrix}

\begin{table}[htbp]
\centering
\caption{Model prediction versus human binary reference label}
\small
\begin{tabular}{@{}lrrr@{}}
\toprule
Human reference & Model operational & Model nonoperational & Weighted total \\
\midrule
Operational & 2,458.69 & 1,885.85 & 4,344.55 \\
Nonoperational & 590.31 & 21,344.15 & 21,934.45 \\
\midrule
Weighted total & 3,049.00 & 23,230.00 & 26,279.00 \\
\bottomrule
\end{tabular}
\end{table}

Weighted performance implied by this matrix is 80.64\% precision, 56.59\%
recall, 97.31\% specificity, and 90.58\% accuracy. Values are population-weighted
estimated counts and need not be integers.

\section{Reproducibility files}

The stable-lexicon and firm-year calculations are generated by
\texttt{analysis\_code/stable\_lexicon\_firm\_check.py}. Machine-readable
outputs are \texttt{reviewer\_final/stable\_lexicon\_check.csv} and
\texttt{reviewer\_final/firm\_level\_proxy\_row.csv}.

\end{document}
